% ============================================================================
%  主文件：main.tex
%  职责：仅做「编排 + 全局开关」，不包含任何正文内容
%  编译：XeLaTeX（latexmk -xelatex，见 .latexmkrc / Makefile）
%  约定：所有正文写在 body/ 与 back/，所有宏与创新命令写在 preamble/
% ============================================================================

% !TEX TS-program = xelatex
% !TEX encoding = UTF-8
\documentclass[review, 12pt, a4paper, onecolumn]{elsarticle}

% ---------------------------------------------------------------------------
% 全局开关（投稿前只改这里）
% ---------------------------------------------------------------------------
\newif\ifdraft
\drafttrue                % 投稿/定稿时改为 \draftfalse：自动关闭行号与 TODO 标记

% ---------------------------------------------------------------------------
% 导言区：分层加载（顺序敏感，勿随意调整）
% ---------------------------------------------------------------------------
% ============================================================================
%  宏包集合：preamble/packages.tex
%  原则：① 只用 TeX Live 主流发行版默认就有的宏包；
%        ② 一个宏包只加载一次，避免重复 \usepackage；
%        ③ 冲突宏包（如 caption/subcaption 与部分期刊类）集中在此处便于排查。
%  编译引擎：XeLaTeX（中文必需，故不加载 inputenc / fontenc）
% ============================================================================

% --- 工具宏（\AtBeginEnvironment 等）---
\usepackage{etoolbox}

% --- 排版基础 ---
\usepackage{xcolor}               % 颜色（其余宏包依赖它）
\usepackage{microtype}            % 字符 protrusion，改善中文排版灰度
\usepackage{geometry}             % 页边距（layout.tex 中设置）

% --- 数学 ---
\usepackage{amsmath, amssymb}     % 公式与数学符号
\usepackage{amsthm}               % 定理/证明/引理环境
\usepackage{mathtools}            % amsmath 增强：\coloneqq, \DeclarePairedDelimiter 等
\usepackage{bm}                   % 加粗数学符号（向量/矩阵）

% --- 插图与表格 ---
\usepackage{graphicx}             % 插图
\usepackage{subcaption}           % 子图（若与期刊模板冲突，改用 \subfigure 或并列 minipage）
\usepackage{booktabs}             % 三线表
\usepackage{array}                % 列格式控制
\usepackage{multirow}             % 跨行单元格
\usepackage{makecell}             % 单元格内换行 \makecell{}
\usepackage{threeparttable}       % 表注（置于表格宽度内）
\usepackage{longtable}            % 跨页长表（符号表适用）
\usepackage{tabularx}             % 自适应宽度列

% --- 算法伪代码 ---
\usepackage{algorithm}
\usepackage{algpseudocode}

% --- 单位与数值 ---
\usepackage{siunitx}              % \SI{}{\metre\per\second}、\num{}
\sisetup{
  range-phrase = {$\sim$},
  list-final-separator = { 和 },
  output-decimal-marker = {.}
}

% --- 中文支持 ---
\usepackage[UTF8]{ctex}           % XeLaTeX + UTF-8 源码，勿改

% --- 列表 ---
\usepackage{enumitem}
\setlist[enumerate]{labelsep=0.6em, leftmargin=2.2em}
\setlist[itemize]{labelsep=0.6em, leftmargin=2.2em}

% --- 审稿辅助（draft 模式下由 draft-tools.tex 决定是否启用）---
\ifdraft
  \usepackage{lineno}
  \modulolinenumbers[5]
\fi
          % 宏包
% ============================================================================
%  版式设置：preamble/layout.tex
%  职责：页面几何、行距、缩进、标题与图表题注格式、浮动体比例
% ============================================================================

% --- 页面几何 ---
\geometry{top=2.5cm, bottom=2.5cm, left=2.8cm, right=2.8cm, headsep=1.0cm}

% --- 段落 ---
\setlength{\parindent}{2em}          % 中文首行缩进两个字符
\setlength{\parskip}{0pt}            % 段落之间不额外加间距（期刊标准做法）
\renewcommand{\arraystretch}{1.2}    % 表格行高，避免内容挤在一起

% --- 浮动体控制（防止图表大量漂移到文末）---
\renewcommand{\topfraction}{0.85}
\renewcommand{\bottomfraction}{0.75}
\renewcommand{\textfraction}{0.10}
\renewcommand{\floatpagefraction}{0.80}
\setcounter{topnumber}{3}
\setcounter{bottomnumber}{2}
\setcounter{totalnumber}{5}

% --- 题注格式 ---
\captionsetup{
  font        = small,
  labelfont   = bf,
  labelsep    = quad,
  figurename  = 图,
  tablename   = 表,
  justification = centering,
  width       = 0.92\linewidth
}

% --- 定理类环境 ---
\newtheoremstyle{plaincn}
  {1.0em}{1.0em}{\itshape}{}{\bfseries}{}{0.6em}{}
\theoremstyle{plaincn}
\newtheorem{theorem}{定理}[section]
\newtheorem{lemma}[theorem]{引理}
\newtheorem{corollary}[theorem]{推论}
\newtheorem{proposition}[theorem]{命题}
\newtheorem{definition}[theorem]{定义}
\newtheorem{assumption}[theorem]{假设}
\newtheorem{remark}[theorem]{注}
            % 版式、间距、中文字号、标题格式
% ============================================================================
%  记号系统：preamble/math-notation.tex
%  ★ 全稿所有数学符号的唯一定义处。新增符号请写在这里并在附录符号表同步。
%  ★ 采用「\x 系列」命名，避免 \a \b \c 之类与已有宏冲突。
% ============================================================================

% --- 常用集合与空间 ---
\newcommand{\Rset}{\mathbb{R}}              % 实数集
\newcommand{\Nset}{\mathbb{N}}              % 自然数集
\newcommand{\Zset}{\mathbb{Z}}              % 整数集

% --- 运算符 ---
\DeclareMathOperator*{\argmin}{arg\,min}    % 带下划线的 argmin
\DeclareMathOperator*{\argmax}{arg\,max}
\DeclareMathOperator{\sign}{sgn}
\DeclareMathOperator{\tr}{tr}               % 矩阵迹

% --- 范数与括号 ---
\DeclarePairedDelimiter{\abs}{\lvert}{\rvert}
\DeclarePairedDelimiter{\norm}{\lVert}{\rVert}
\DeclarePairedDelimiter{\inner}{\langle}{\rangle}
\DeclarePairedDelimiterX{\set}[1]{\{}{\}}{#1}

% --- 偏导与微分 ---
\newcommand{\pd}[2]{\frac{\partial #1}{\partial #2}}
\newcommand{\dd}[1]{\mathrm{d}#1}

% --- 论文专属符号（示例，按实际模型增删）---
% 风险类
\newcommand{\pcrash}[0]{P_{\mathrm{crash}}}              % 坠机概率
\newcommand{\pcrashT}[2]{\pcrash\!\left(#1 \mid #2\right)} % 时变坠机概率 P(t | x, \theta)
\newcommand{\hazard}[0]{\lambda}                          % 危险率
\newcommand{\survival}[0]{S}                              % 生存函数
\newcommand{\exposure}[0]{\rho}                           % 暴露密度
\newcommand{\casualty}[0]{C_{\mathrm{fat}}}               % 致死代价
\newcommand{\propertyloss}[0]{C_{\mathrm{prop}}}          % 财产代价
\newcommand{\noisecost}[0]{C_{\mathrm{noise}}}            % 噪声代价
\newcommand{\riskcost}[0]{\mathcal{R}}                    % 综合风险代价
\newcommand{\riskfield}[0]{\mathcal{F}}                   % 风险张量场

% 规划类
\newcommand{\statespace}[0]{\mathcal{S}}    % 状态空间
\newcommand{\action}[0]{\mathcal{A}}        % 动作集
\newcommand{\nodeset}[0]{\mathcal{V}}       % 节点集
\newcommand{\edgeset}[0]{\mathcal{E}}       % 边集
\newcommand{\heur}[0]{h}                    % 启发函数
\newcommand{\gcost}[0]{g}                   % 累积代价
\newcommand{\tlabel}[0]{\tau}               % 时间标签

% --- 矩阵转置与向量 ---
\newcommand{\transpose}[0]{^{\mathsf{T}}}
\newcommand{\superscript}[1]{^{\mathrm{#1}}}
\newcommand{\subscript}[1]{_{\mathrm{#1}}}
     % 数学符号与记号（全稿唯一定义处）
% ============================================================================
%  文本层宏：preamble/text-macros.tex
%  职责：术语缩写、中英混排、作者自定义短语
%  原则：凡是「可能全文统一替换」的写法，一律定义成宏（改一次全局生效）
% ============================================================================

% --- 拉丁文缩写 ---
\newcommand{\etal}{等}
\newcommand{\ie}{即}
\newcommand{\eg}{例如}
\newcommand{\etc}{等}
\newcommand{\wrt}{相对于}
\newcommand{\st}{使得}

% --- 常用中英术语（保证全文写法统一）---
\newcommand{\UAV}{UAV}
\newcommand{\TPR}{TPR}
\newcommand{\termTPR}{第三方风险（third-party risk, \TPR）}
\newcommand{\termTD}{\mbox{时间依赖}}
\newcommand{\methodname}{TD-RiskA\textsuperscript{*}}   % 本文算法名称
\newcommand{\methodnamelatex}{TD-RiskA$^*$}
\newcommand{\basestatic}{Static-RiskA$^*$}              % 基线方法示例

% --- 作者信息快捷引用 ---
\newcommand{\Corresponding}{通讯作者}
\newcommand{\CoFirst}{共同第一作者}

% --- 强调与浮动对象 ---
\newcommand{\keywordcn}[1]{\textbf{#1}}
\newcommand{\file}[1]{\texttt{#1}}
\newcommand{\code}[1]{\texttt{#1}}
       % 文本层自定义命令（术语、缩写、单位）
% ============================================================================
%  伪代码环境：preamble/algorithms.tex
%  约定：仅重定义 algpseudocode 中确定存在的关键字，避免因版本差异报错。
%        正文中使用 \begin{algorithmic} 块，\Input / \Output 已中文化。
% ============================================================================

\algnewcommand\algorithmicinput{\textbf{输入：}}
\algnewcommand\algorithmicoutput{\textbf{输出：}}
\algnewcommand\Input{\item[\algorithmicinput]}
\algnewcommand\Output{\item[\algorithmicoutput]}

\algrenewcommand\algorithmicrequire{\textbf{输入：}}
\algrenewcommand\algorithmicensure{\textbf{输出：}}

% 伪代码体统一用脚注字号 + 单栏，编号按节（第 5 节 → 算法 5.1）
\renewcommand{\thealgorithm}{\arabic{section}.\arabic{algorithm}}
\AtBeginEnvironment{algorithm}{\footnotesize}
        % 伪代码环境中文配置
% ============================================================================
%  写作辅助：preamble/draft-tools.tex
%  仅在草稿模式（\drafttrue）下显示批注与占位图框；
%  切成 \draftfalse 后，一切自动消失，直接产出可投稿版本。
%  依赖：xcolor（已在 packages.tex 加载）—— 不引入额外宏包，降低环境风险。
%  注：批注一律用「行内」而非页边 \marginpar，
%      因为 marginpar 在 frontmatter / 标题框 / 表格内会触发 Float(s) lost。
% ============================================================================

\definecolor{todocolor}{RGB}{200, 40, 40}
\definecolor{notecolor}{RGB}{40, 90, 180}
\definecolor{phcolor}{RGB}{140, 140, 140}

\ifdraft

  % TODO：红色行内标记，编译后可在 PDF 里直接检索 "[TODO"
  \newcommand{\todo}[1]{\textcolor{todocolor}{\textbf{[TODO]~#1}}}

  % 作者批注：蓝色行内标记，可带作者名：\note[张三]{建议再补一组消融}
  \newcommand{\note}[2][note]{\textcolor{notecolor}{$\langle$#1：#2$\rangle$}}

  % 占位图：尚未出图的灰色占位框，#1 = 宽度比例，#2 = 占位说明
  \newcommand{\placeholderfig}[2][0.8]{%
    \fbox{\parbox{\dimexpr#1\linewidth-2\fboxsep-2\fboxrule\relax}{%
      \centering\vspace{2.2cm}%
      \textcolor{phcolor}{(占位图) #2}%
      \vspace{2.2cm}}}%
  }

\else

  % 定稿模式：批注全部隐形（正文中的命令无需逐个删除）
  \newcommand{\todo}[1]{\relax}
  \newcommand{\note}[2][note]{\relax}
  \newcommand{\placeholderfig}[2][0.8]{\textcolor{phcolor}{(待补插图：#2)}}

\fi
       % TODO / 批注 / 占位图工具
% ============================================================================
%  引用与链接：preamble/references-setup.tex
%  ★ 必须最后加载：hyperref 之后再加载 cleveref，否则交叉引用名称会错乱。
% ============================================================================

\definecolor{linkcolor}{RGB}{0, 102, 204}
\definecolor{citecolor}{RGB}{0, 128, 0}
\definecolor{urlcolor}{RGB}{85, 26, 139}

\usepackage{hyperref}
\hypersetup{
  colorlinks     = true,
  linkcolor      = linkcolor,
  citecolor      = citecolor,
  urlcolor       = urlcolor,
  pdfstartview   = {Fit},
  bookmarksnumbered = true,
  bookmarksopen  = true,
}

% pdftitle / pdfauthor 依赖 config/metadata.tex 中的宏，故延后到开文时才展开
\AtBeginDocument{%
  \hypersetup{pdftitle={\papertitle}, pdfauthor={\authorname}}%
}

\usepackage[nameinlink]{cleveref}

% --- cleveref 中文名称 ---
\crefname{section}{节}{节}
\Crefname{section}{节}{节}
\crefname{figure}{图}{图}
\Crefname{figure}{图}{图}
\crefname{table}{表}{表}
\Crefname{table}{表}{表}
\crefname{equation}{式}{式}
\Crefname{equation}{式}{式}
\crefname{algorithm}{算法}{算法}
\Crefname{algorithm}{算法}{算法}
\crefname{appendix}{附录}{附录}
\Crefname{appendix}{附录}{附录}
\crefname{theorem}{定理}{定理}
\crefname{lemma}{引理}{引理}
\crefname{definition}{定义}{定义}

% --- 语义化引用快捷命令（推荐优先于裸 \ref）---
% 注意：不要取 \reffig/\reftab/\refalg 等名字，它们可能与 cleveref 自动生成的
%       \<float>ref 系列冲突（实测 \refalg 会被 cleveref 占用）。统一用 \ref* 前缀。
\newcommand{\reffig}[1]{图~\ref{#1}}
\newcommand{\reftab}[1]{表~\ref{#1}}
\newcommand{\refsec}[1]{第~\ref{#1}~节}
\newcommand{\refeqn}[1]{式~(\ref{#1})}
\newcommand{\refalg}[1]{算法~\ref{#1}}
\newcommand{\refapp}[1]{附录~\ref{#1}}

% 说明：elsarticle 采用数字制引文，直接使用 \cite{} 即可（多文献自动压缩为 [1,3--5]）。
  % hyperref 与 cleveref（必须最后加载）

% ---------------------------------------------------------------------------
% 元数据：标题、作者、摘要、关键词、亮点
% ---------------------------------------------------------------------------
% ============================================================================
%  元数据：config/metadata.tex
%  ★ 所有「会随投稿/改稿而变」的元信息集中在此，正文与导言区均不含硬编码。
% ============================================================================

% --- 标题 ---
% 建议：<主题>：<方法>与<验证>，控制在 15 个词以内
\newcommand{\papertitle}{城市低空无人机风险感知航路规划：时空异质性建模与自适应求解}

% --- 作者简介 ---
\newcommand{\authorname}{作者一，作者二，作者三}
\newcommand{\authoremail}{author1@example.com}
\newcommand{\correspondingnote}{通讯作者}
\newcommand{\authoraddress}{某某大学某某学院，城市 邮编}

% --- 基金 ---
\newcommand{\fundingname}{国家自然科学基金}
\newcommand{\fundingnumber}{No. 00000000}

% --- MSC 分类码（若目标期刊不需要，留空即可）---
\newcommand{\msccodes}{90B20 \sep 90C59 \sep 68T40}

% ============================================================================
%  摘要与关键词：config/abstract.tex
%  写法建议：背景(1) → 缺口(1) → 方法(2) → 结果(1) → 意义(1)
% ============================================================================

\newcommand{\abstracttext}{%
【背景】无人机在城市低空的规模化运行，使其失效坠落对地面第三方造成的风险成为监管与运营的核心约束。\quad
【缺口】既有风险感知路径规划研究多假定坠落概率与人群暴露恒定的静态风险场，无法刻画微气象、建筑峡谷效应与人口昼夜潮汐共同导致的时空异质性。\quad
【方法】针对这一问题，本文构建了……

\todo{按「背景—缺口—方法—结果—意义」五句法补全摘要，中英文各一版，控制在 200～250 字。}
}

% --- 关键词：5～6 个，首字母无需大写 ---
\newcommand{\paperkeywords}{%
第三方风险 \sep 时空异质性 \sep 无人机 \sep 低空空域 \sep 路径规划 \sep 风险-成本权衡
}

% ============================================================================
%  研究亮点：config/highlights.tex
%  Elsevier 要求 3～5 条，每条不超过 85 个字符，动词开头、结论式表述。
% ============================================================================

\newcommand{\researchhighlights}{%
\item 提出面向……的时变第三方风险张量模型，将气象、地表与人口节律耦合进统一代价测度。
\item 设计……算法，在不膨胀搜索图的前提下实现时间依赖代价的最优路径求解。
\item 证明在单调代价结构下……，为启发式剪枝提供理论依据。
\item 在合成与真实城市场景中验证：相比最优基线，……指标下降……\%，绕行代价低于……\%。
}


% ---------------------------------------------------------------------------
% 资源路径
% ---------------------------------------------------------------------------
\graphicspath{{assets/figures/}}       % 插入插图时只需写文件名，不带路径

% ===========================================================================
\begin{document}

% ---------------------------------------------------------------------------
% 前置部分
% ---------------------------------------------------------------------------
\ifdraft\linenumbers\fi                % 审稿行号（draft 模式下每 5 行一个编号）

\begin{frontmatter}

\title{\papertitle}

\author[addr1]{\authorname\corref{cor1}\fnref{fn1}}
\ead{\authoremail}
\cortext[cor1]{\correspondingnote}
\fntext[fn1]{本工作得到了 \fundingname 的资助（项目编号 \fundingnumber）。}
\address[addr1]{\authoraddress}

\begin{abstract}
\abstracttext
\end{abstract}

% 图形摘要（可选，按需取消注释）：一张图概括全文研究逻辑。
% 注意：留空（仅保留环境但不放内容）会触发 “Float(s) lost” 错误，故默认注释。
% \begin{graphicalabstract}
%   \includegraphics[width=\linewidth]{graphical-abstract.pdf}
% \end{graphicalabstract}

\begin{highlights}
  \researchhighlights
\end{highlights}

\begin{keyword}
\paperkeywords
\MSC[2020]{\msccodes}
\end{keyword}

\end{frontmatter}

% ---------------------------------------------------------------------------
% 正文：按章节编号前缀顺序加载（body/NN-name.tex）
% ---------------------------------------------------------------------------
% ============================================================================
%  第 1 章 引言 —— body/01-introduction.tex
%  职责：讲清「为什么必须做」与「本文做了什么」，行文以问题驱动。
% ============================================================================

\section{引言}
\label{sec:introduction}

% ---------------------------------------------------------------------------
\subsection{研究背景}

无人机（\UAV）在城市低空的物流配送、设施巡检与应急响应中的应用规模持续增长，
随之而来的是其对地面人群、车辆与建筑物构成的 \termTPR。

\todo{补一到两句宏观数据：低空经济规模、飞行架次增速，注意引用权威统计源。}

% ---------------------------------------------------------------------------
\subsection{问题提出}

既有空域审批与航线规划多基于飞行前的静态快照，而实际运行环境随
时间尺度剧烈变化：

\begin{itemize}
  \item \textbf{微气象时变}：风速、风向与降雨强度改变失效概率；
  \item \textbf{建筑峡谷扰动}：街谷内的湍流与遮挡放大局部不确定性；
  \item \textbf{人口潮汐}：同一位置在通勤高峰与深夜的暴露人群差异达数量级。
\end{itemize}

由此形成「静态审批—动态环境」之间的安全裂隙，构成本文的核心问题：
\textit{如何在风险随时空变化的条件下，为给定出发时刻规划一条长期最优且在线可行的路径？}

% ---------------------------------------------------------------------------
\subsection{现有研究的局限}

从三个维度评述既有工作的不足（详见 \refsec{sec:related-work}）：
模型维度单一、时间维度缺失、求解维度保守。

% ---------------------------------------------------------------------------
\subsection{主要贡献}

本文的贡献可概括为四点：

\begin{enumerate}
  \item \textbf{模型}：构建……
  \item \textbf{算法}：提出 \methodnamelatex，……
  \item \textbf{理论}：证明……
  \item \textbf{实证}：在……场景验证……
\end{enumerate}

\todo{贡献清单要与 \refsec{sec:related-work} 的表格 Columns 一栏严丝合缝地对应。}

% ---------------------------------------------------------------------------
\subsection{文章结构}

全文组织如下：\refsec{sec:related-work} 回顾相关工作；
\refsec{sec:preliminaries} 给出问题定义；
\refsec{sec:risk-model} 与 \refsec{sec:planning} 分别阐述模型与算法；
\refsec{sec:experiments} 给出实验；\refsec{sec:discussion} 讨论；
\refsec{sec:conclusion} 总结。

% ============================================================================
%  第 2 章 相关工作 —— body/02-related-work.tex
%  职责：把本文放回到坐标系里，突出「别人没做的那格」。
%  写法：分类评述 → 对比表 → 明确空白
% ============================================================================

\section{相关工作}
\label{sec:related-work}

% ---------------------------------------------------------------------------
\subsection{第三方风险建模}

坠机概率模型经历了从常数假设向动态建模的演进。早期研究以恒定失效率
或经验图谱刻画 $\pcrash$ \citep{primatesta2022review}，
难以反映环境状态对失效概率的影响。近期工作引入比例风险模型处理
微气象协变量 \citep{cooper2023proportional}，但在暴露度与时间维度上
仍沿用静态假设。

% ---------------------------------------------------------------------------
\subsection{风险感知路径规划}

基于风险场的路径搜索方法可归纳为两类：栅格/概率地图上的改进启发式搜索
\citep{zhang2024spatiotemporal}，以及以可达安全性为约束的优化/采样方法
\citep{chen2024costastar}。前者多采用静态风险场，后者则通常假设
滑动窗内的短期最优。

% ---------------------------------------------------------------------------
\subsection{时间与时空依赖的搜索方法}

当边代价随通过时刻变化时，问题退化为时间依赖最短路：代价函数须满足
FIFO（先到先走）性质，标签更改/标签更正算法构成主流求解框架。
\todo{补 citation：time-dependent shortest path 经典文献。}

% ---------------------------------------------------------------------------
\subsection{研究空白定位}

\reftab{tab:literature-comparison} 从「风险维度、时间动态、暴露度演化、求解方式」
四个维度对比代表性工作，据此定位本文所处的空白。

% 表格统一放在 assets/tables/ 并以 \input 引入，便于多人协作时分离文件冲突
% ============================================================================
%  文献对比表 —— assets/tables/literature-comparison.tex
%  位置：第 \ref{sec:related-work} 节
% ============================================================================

\begin{table}[t]
  \centering
  \caption{相关工作在四个维度上的对比（$\checkmark$~支持，$	imes$~不支持，$\sim$~部分支持）}
  \label{tab:literature-comparison}
  \footnotesize
  \begin{threeparttable}
  \begin{tabular}{@{}lcccc@{}}
    \toprule
    工作 & 多维风险 & 时间动态 & 暴露度演化 & 最优性保证 \\
    \midrule
    Primatesta 等 \cite{primatesta2022review} & $\checkmark$        & $	imes$           & $	imes$           & --          \\
    Cooper 等 \cite{cooper2023proportional}   & $	imes$            & $\checkmark$       & $	imes$           & --          \\
    Zhang 等 \cite{zhang2024spatiotemporal}   & $\sim$              & $\sim$             & $	imes$           & $	imes$    \\
    Chen 等 \cite{chen2024costastar}          & $	imes$            & $\sim$             & $	imes$           & $\sim$      \\
    \midrule
    \textbf{本文} & $\checkmark$ & $\checkmark$ & $\checkmark$ & $\checkmark$ \\
    \bottomrule
  \end{tabular}
  \begin{tablenotes}
    \footnotesize
    \item[] 注：$\checkmark$/$	imes$/$\sim$ 分别表示支持、不支持、部分支持。
  \end{tablenotes}
  \end{threeparttable}
\end{table}


% ============================================================================
%  第 3 章 预备知识与问题定义 —— body/03-preliminaries.tex
%  职责：符号、空域离散化、时空状态空间、优化目标。
% ============================================================================

\section{预备知识与问题定义}
\label{sec:preliminaries}

% ---------------------------------------------------------------------------
\subsection{空域离散化}

将城市低空空域离散为三维栅格节点集 $\nodeset$，记节点 $v \in \nodeset$ 的
空间坐标为 $\mathbf{x}_v = (x_v, y_v, z_v)^{\top}$，
可行转移边集为 $\edgeset$。

\todo{补网格分辨率的取值依据与敏感性分析指向（→ \refsec{sec:experiments}）。}

% ---------------------------------------------------------------------------
\subsection{时空状态空间}

引入离散时间索引 $t \in \mathcal{T}$，定义时空状态
\begin{equation}
  s = (v, t) \in \statespace = \nodeset \times \mathcal{T},
  \quad |\mathcal{T}| = T.
  \label{eq:state-def}
\end{equation}

符号表见 \refapp{app:notation}，全文符号均在
\file{preamble/math-notation.tex} 中统一定义。

% ---------------------------------------------------------------------------
\subsection{优化目标}

对给定起点 $v_0$、终点 $v_d$ 与出发时刻 $t_0$，求解
\begin{subequations}
\label{eq:objective}
\begin{alignat}{2}
  & \min_{\pi \in \Pi}        & \quad & J(\pi) = \sum_{k=0}^{K-1}
      \riskcost\bigl((v_k, t_k) \to (v_{k+1}, t_{k+1})\bigr) \\
  & \text{s.t.}               &       & (v_k, t_k) \in \statespace,\;
      t_{k+1} = t_k + \Delta t(v_k, v_{k+1}), \\
  &                           &       & v_K = v_d,\; t_K \le t_0 + T_{\max},
\end{alignat}
\end{subequations}
其中 $\riskcost$ 为 \refsec{sec:risk-model} 定义的综合风险代价，$\Delta t$ 为
边长对应的飞行时间。

\begin{assumption}[代价可加性]
综合风险代价沿路径可加，即
$J(\pi) = \sum_k \riskcost_k$，且 $\riskcost \ge 0$。
\end{assumption}

% ============================================================================
%  第 4 章 风险建模 —— body/04-risk-model.tex
%  职责：把「风险的成本」拆成人能验证的子模块：概率 × 暴露 × 后果 × 集成
% ============================================================================

\section{时空异质性风险建模}
\label{sec:risk-model}

% ---------------------------------------------------------------------------
\subsection{时变失效概率}

沿用比例风险模型的构造，令基线危险率为 $\hazard_0(t)$，协变量向量
$\mathbf{z}(t)$ 包含风速、降雨与建筑峡谷因子，则
\begin{equation}
  \pcrashT{t}{\mathbf{x}, \mathbf{z}}
    = 1 - \exp\!\left[-\int_{0}^{t} \hazard_0(u)
      \exp\bigl(\boldsymbol{\beta}\transpose \mathbf{z}(u)\bigr) \, \dd{u}\right],
  \label{eq:hazard}
\end{equation}
其中 $\boldsymbol{\beta}$ 为回归系数向量，其物理意义为各协变量对
失效风险的弹性。

\todo{说明 $\hazard_0$ 的分布族选择（Weibull / 常数）与标定数据来源。}

% ---------------------------------------------------------------------------
\subsection{暴露度演化}

地表 $h$ 处单元 $m$ 在时刻 $t$ 的暴露密度由静态基底与活动强度合成：
\begin{equation}
  \exposure(m, t) = \exposure_{\mathrm{static}}(m)
    + \exposure_{\mathrm{tidal}}(m) \, \phi\bigl(t; \pi(m)\bigr),
  \label{eq:exposure}
\end{equation}
其中 $\pi(m)$ 标识土地利用类型，$\phi(\cdot)$ 为对应的昼夜激活曲线。

% ---------------------------------------------------------------------------
\subsection{后果代价：致死、财产与噪声}

三个维度的后果代价分别为 $\casualty$、$\propertyloss$ 与 $\noisecost$。
以致死代价为例：
\begin{equation}
  \casualty(m, t) = \exposure(m, t) \, A_{\mathrm{exp}}
    \, \pcrashT{t}{\mathbf{x}_m, \mathbf{z}}.
  \label{eq:casualty}
\end{equation}

% ---------------------------------------------------------------------------
\subsection{综合集成与性质}

\begin{definition}[综合风险代价]
给定边 $(u \to v)$ 与其通过时刻 $t$，定义
\begin{equation}
  \riskcost\bigl((u, t) \to (v, t')\bigr)
    = w_1 \casualty + w_2 \propertyloss + w_3 \noisecost,
  \quad w_i \ge 0,\; \sum_i w_i = 1.
  \label{eq:risk-cost}
\end{equation}
\end{definition}

\begin{theorem}[单调性]
\label{thm:monotone}
若 $\hazard_0(t)$ 与 $\boldsymbol{\beta}\transpose \mathbf{z}(t)$ 均非负有界，
则由 \refeqn{eq:risk-cost} 定义的边代价满足
$\riskcost(t_2) \ge \riskcost(t_1)$ 对 $t_2 > t_1$ 未必成立，
但其积分形式对时间可测且满足 Bellman 最优性方程的前提条件。
\end{theorem}

\begin{proof}
证明见 \refapp{app:supplement}。
\todo{补完整证明：核心是证明可加性 + FIFO + 非负性三条。}
\end{proof}

% ============================================================================
%  第 5 章 求解算法 —— body/05-planning-algorithm.tex
%  职责：给出算法结构（含伪代码）、关键机制、复杂度与最优性。
% ============================================================================

\section{\termTD 风险感知路径规划算法}
\label{sec:planning}

% ---------------------------------------------------------------------------
\subsection{总体思路}

直接在三维节点上再展开一维时间，会把搜索图规模放大 $T$ 倍。
\methodnamelatex 的核心是把时间从\textbf{图拓扑}中剥离，作为
\textbf{标签状态}沿路径传播，从而在不膨胀搜索图的前提下支持任意时刻的
在线风险查询。

% ---------------------------------------------------------------------------
\subsection{标签化传播}

对节点 $v$，维护标签 $\ell = (\gcost, \tlabel, \mathrm{parent})$，
其中 $\tlabel$ 为到达时刻；松弛时按
\refeqn{eq:state-def} 更新时间：
\begin{equation}
  \tlabel(v') = \tlabel(v) + \Delta t(v, v'),
  \qquad
  \gcost(v') = \gcost(v) + \riskcost\bigl((v, \tlabel(v)) \to (v', \tlabel(v'))\bigr).
  \label{eq:label-prop}
\end{equation}

% ---------------------------------------------------------------------------
\subsection{累积危险率的加法形式}

为避免长路径上概率乘积导致的数值下溢，改以累积危险率传播：
\begin{equation}
  H(v') = H(v) + \int_{\tlabel(v)}^{\tlabel(v')} \hazard(u) \, \dd{u},
  \qquad
  \survival = \exp(-H).
  \label{eq:cumulative-hazard}
\end{equation}

% ---------------------------------------------------------------------------
\subsection{算法流程}

\begin{algorithm}[t]
\caption{\methodnamelatex：时间依赖风险感知最短路搜索}
\label{alg:td-risk-astar}
\begin{algorithmic}[1]
  \Input 时空图 $G = (\statespace, \edgeset)$，起点 $(v_0, t_0)$，终点 $v_d$，权重 $\{w_i\}$
  \Output 最优路径 $\pi^{*}$
  \State $\mathrm{OPEN} \gets \{(v_0, \gcost{=}0, \tlabel{=}t_0)\}$，$\mathrm{CLOSED} \gets \emptyset$
  \While {$\mathrm{OPEN} \neq \emptyset$}
    \State $\ell \gets \argmin\{\gcost + \heur(\cdot) \mid \ell \in \mathrm{OPEN}\}$
    \State 从 $\mathrm{OPEN}$ 取出 $\ell$
    \If {$\ell.v = v_d$}
      \State \textbf{return} 回溯标签链得到 $\pi^{*}$
    \EndIf
    \For {每个后继 $v' \in \mathcal{N}(\ell.v)$}
      \State $\tlabel' \gets \ell.\tlabel + \Delta t(\ell.v, v')$
      \State $\riskcost' \gets \textsc{RiskQuery}(\ell.v, v', \tlabel')$       \Comment 在线查询动态风险张量
      \State $\ell' \gets (\gcost{=}\ell.\gcost + \riskcost',\; \tlabel{=}\tlabel',\; \mathrm{parent}{=}\ell)$
      \If {$\ell'$ 被 $\mathrm{CLOSED}(v')$ 中的标签支配} \Comment 见第 \ref{sec:dominance} 节
        \State \textbf{continue}
      \EndIf
      \State 更新 $\mathrm{OPEN}(v')$ 与 $\mathrm{CLOSED}(v')$
    \EndFor
  \EndWhile
\end{algorithmic}
\end{algorithm}

% ---------------------------------------------------------------------------
\subsection{支配剪枝}
\label{sec:dominance}

\begin{definition}[标签支配]
标签 $\ell_1$ 支配 $\ell_2$，当且仅当
$\gcost_1 \le \gcost_2$ 且 $\tlabel_1 \le \tlabel_2$，且至少一处严格不等。
\end{definition}

\begin{theorem}[剪枝安全性（单标签覆盖）]
在 \cref{thm:monotone} 的单调代价结构下，被支配标签不可能导出优于保留标签的解。
\end{theorem}

% \todo{补充剪枝正确性的完整证明，并与 \cref{thm:monotone} 的前提条件对齐。}

% ---------------------------------------------------------------------------
\subsection{复杂度分析}

设平均出度为 $b$、路径长度为 $K$、支配检查代价为 $O(\log q)$（$q$ 为标签数上界），
则时间复杂度为 $O(bK \log q)$，空间复杂度为 $O(|\nodeset| \, q)$；
在 $q$ 有界的情形下与传统三维 \basestatic \ 同阶。

% ============================================================================
%  第 6 章 实验 —— body/06-experiments.tex
%  职责：可复现。设置 → 基线 → 结果（图/表）→ 消融 → 效率。
% ============================================================================

\section{实验与结果}
\label{sec:experiments}

% ---------------------------------------------------------------------------
\subsection{实验设置}

\begin{table}[t]
  \centering
  \caption{主要实验参数}
  \label{tab:settings}
  \begin{threeparttable}
  \begin{tabular}{@{}llcc@{}}
    \toprule
    类别 & 参数 & 取值 & 单位 \\
    \midrule
    空域   & 栅格分辨率     & \num{50}            & \si{\metre} \\
    时间   & 时间步长       & \num{60}            & \si{\second} \\
    风险   & 权重 $w_1$/$w_2$/$w_3$ & 0.5/0.2/0.3 & -- \\
    算法   & 标签上限 $q$   & \num{3}             & -- \\
    \bottomrule
  \end{tabular}
  \begin{tablenotes}
    \footnotesize \item[] 注：权重组合的敏感性分析见 \refsec{sec:sensitivity}。
  \end{tablenotes}
  \end{threeparttable}
\end{table}

基线方法选取 \basestatic（静态风险场）与 Greedy-Risk（贪心）。

% ---------------------------------------------------------------------------
\subsection{合成场景}

\begin{figure}[t]
  \centering
  \begin{subfigure}[b]{0.48\linewidth}
    \centering
    \placeholderfig{静态风险场下的规划结果}
    \caption{静态风险场}
    \label{fig:synthetic-static}
  \end{subfigure}
  \hfill
  \begin{subfigure}[b]{0.48\linewidth}
    \centering
    \placeholderfig{本文方法结果}
    \caption{\methodnamelatex（本文）}
    \label{fig:synthetic-ours}
  \end{subfigure}
  \caption{合成场景中不同时刻的规划路径对比。}
  \label{fig:synthetic}
\end{figure}

如 \reffig{fig:synthetic} 所示，……\todo{补 2～3 句结果分析，务必给出定量指标。}

% ---------------------------------------------------------------------------
\subsection{真实城市场景}

% ============================================================================
%  主结果表 —— assets/tables/main-results.tex
%  位置：第 \ref{sec:experiments} 节
% ============================================================================

\begin{table}[t]
  \centering
  \caption{真实城市场景下的主要结果（均值 $\pm$ 标准差；风险类指标越低越好，路径长度自然越低越好）}
  \label{tab:main-results}
  \begin{threeparttable}
  \begin{tabular}{@{}lcccc@{}}
    \toprule
    方法 & 综合风险 & 致死风险 & 噪声影响 & 路径长度 (km) \\
    \midrule
    \basestatic     & xx.xx $\pm$ x.xx & xx.xx & xx.xx & xx.xx \\
    Greedy-Risk     & xx.xx $\pm$ x.xx & xx.xx & xx.xx & xx.xx \\
    \textbf{本文}   & \textbf{xx.xx} $\pm$ x.xx & \textbf{xx.xx} & xx.xx & xx.xx \\
    \midrule
    相对最优基线   & $-$\,xx.x\% & $-$\,xx.x\% & $-$\,xx.x\% & $+\,$x.x\% \\
    \bottomrule
  \end{tabular}
  \begin{tablenotes}
    \footnotesize
    \item[] 注：最优值加粗显示；相对值为各方法的均值对比。
  \end{tablenotes}
  \end{threeparttable}
\end{table}


% ---------------------------------------------------------------------------
\subsection{消融实验}

依次移除时间维传播、累积危险率形式与支配剪枝，观察各项指标变化，
验证各模块贡献。

% ---------------------------------------------------------------------------
\subsection{计算效率}

记录各方法在相同地图上的运行时间与扩展节点数：在不显著增加运行时间的
前提下，本文方法将……$X\%$。

% ---------------------------------------------------------------------------
\subsection{敏感性分析}
\label{sec:sensitivity}

考察权重 $w_i$、栅格分辨率与时间步长对结果的影响。

% ============================================================================
%  第 7 章 讨论 —— body/07-discussion.tex
%  职责：解释为什么结果成立、在哪有效、在哪失效。避免重复「结果」。
% ============================================================================

\section{讨论}
\label{sec:discussion}

% ---------------------------------------------------------------------------
\subsection{结果的解释}

本文方法之所以在昼夜节律显著的区域收益最大，原因在于：
风险在同一地点的跨时刻差异被显式建模，迁航决策因而获得
「择时绕行」的自由度。\todo{把「为什么」讲透，而不是重复「降了多少」。}

% ---------------------------------------------------------------------------
\subsection{适用边界}

在气象条件平稳、人口分布均匀的城市外围，模型增益有限；此时
选用更简单的静态场即可，不必承担额外的在线查询开销。

% ---------------------------------------------------------------------------
\subsection{有效性威胁}

\begin{itemize}
  \item \textbf{构念效度}：以……作为风险代理指标，可能遗漏……；
  \item \textbf{内部效度}：参数标定依赖……数据集，存在标定误差；
  \item \textbf{外部效度}：实验限于……城市，向其他形态城市外推需谨慎。
\end{itemize}

% ============================================================================
%  第 8 章 结论 —— body/08-conclusion.tex
%  写法：① 一句话总结 ② 三点贡献（不复述实验数字）③ 局限与未来工作
% ============================================================================

\section{结论}
\label{sec:conclusion}

本文围绕城市低空无人机运行中的第三方风险时空异质性问题，建立了……
并以……求解，理论与时实验两方面验证了……

具体而言：
\begin{enumerate}
  \item 在建模上，……；
  \item 在算法上，……；
  \item 在实证上，……。
\end{enumerate}

本文的局限在于……；后续研究将考虑多机协同场景下的风险耦合建模、
以及面向在线部署的增量式风险张量更新机制。

\todo{结论不要出现新的引用与新术语，只做收束。}


% ---------------------------------------------------------------------------
% 尾部：声明、致谢、附录
% ---------------------------------------------------------------------------
% ============================================================================
%  期刊声明 —— back/declarations.tex
%  Elsevier 等出版商要求的标准声明；投稿前按目标期刊口径逐项修改。
%  若目标期刊不需要某一项，直接删除对应 \section* 即可。
% ============================================================================

\section*{作者贡献声明（CRediT）}

\textbf{作者一}：概念化，方法论，形式化分析，初稿撰写。
\textbf{作者二}：软件实现，数据整理，可视化。
\textbf{作者三}：指导，审阅与修订，资金获取。

\section*{资金来源声明}

本工作得到了 \fundingname（项目批准号 \fundingnumber）的资助。
资助方未参与研究设计、数据采集与论文撰写过程。

\section*{利益冲突声明}

作者声明不存在任何可能影响本研究的竞争性财务利益或个人关系。

\section*{数据与代码可用性声明}

本研究使用的数据与代码可在 \url{https://github.com/xxx/xxx} 获取；
涉密/受限数据依据相关协议向通信作者申请获得。

\section*{生成式 AI 使用声明}

在本文撰写过程中，作者使用了 XXX 工具进行语言润色与代码辅助，
所有内容经作者审阅并对文稿的准确性与完整性负责。

\section*{伦理声明}

本研究不涉及人体或动物实验。

% ============================================================================
%  致谢 —— back/acknowledgments.tex
% ============================================================================

\section*{致谢}

感谢 XXX 提供的计算资源与数据支持；感谢 XXX 在方法讨论中的建设性意见。


\appendix
% ============================================================================
%  附录 A：符号表 —— back/appendix-a-notation.tex
%  注意：本文件由 main.tex 在 \appendix 之后加载，故编号自动为 A、B……
% ============================================================================

\section{符号表}
\label{app:notation}

% 符号改动必须同步 preamble/math-notation.tex 与本表，避免文稿前后不一致
% ============================================================================
%  符号表 —— assets/tables/notation.tex
%  位置：附录 \ref{app:notation}
%  规范：改动符号 → 同步本表 + preamble/math-notation.tex
% ============================================================================

\small
\renewcommand{\arraystretch}{1.15}
\begin{longtable}{@{}p{0.16\linewidth}p{0.62\linewidth}l@{}}
  \caption{主要符号一览} \label{tab:notation} \\
  \toprule
  符号 & 含义 & 单位 \\
  \midrule
  \endfirsthead
  \multicolumn{3}{l}{\tablename~\thetable（续）} \\
  \toprule
  符号 & 含义 & 单位 \\
  \midrule
  \endhead
  \bottomrule
  \endfoot

  $\nodeset, \edgeset$    & 空域节点集与可行边集 & -- \\
  $\statespace$           & 时空状态空间 & -- \\
  $t, \mathcal{T}$        & 时刻与时间索引集 & \si{\second} \\
  $\pcrash$               & 坠机（失效）概率 & -- \\
  $\hazard, \hazard_0$    & 危险率与基线危险率 & \si{\per\second} \\
  $\survival$             & 生存函数 & -- \\
  $\exposure$             & 地面暴露密度 & \si{\per\square\metre} \\
  $\casualty$             & 致死代价 & 元 \\
  $\propertyloss$         & 财产损失代价 & 元 \\
  $\noisecost$            & 噪声影响代价 & 元 \\
  $\riskcost$             & 综合风险代价 & 元 \\
  $w_1, w_2, w_3$         & 三维风险权重 & -- \\
  $\gcost, \heur$         & 累积代价与启发函数 & 元 \\
  $\tlabel$               & 到达时刻标签 & \si{\second} \\
\end{longtable}
\renewcommand{\arraystretch}{1.2}


% ============================================================================
%  附录 B：证明与补充材料 —— back/appendix-b-supplement.tex
% ============================================================================

\section{证明与补充说明}
\label{app:supplement}

\subsection{\cref{thm:monotone} 的证明}

由 \refeqn{eq:hazard} 出发，结合积分中值定理可得……\todo{补证明步骤。}

\subsection{补充实验结果}

\begin{table}[htbp]
  \small
  \centering
  \caption{补充实验：不同天气等级下的风险节省率}
  \label{tab:supplement-weather}
  \begin{tabular}{@{}lccc@{}}
    \toprule
    天气等级 & 本文 & \basestatic & 节省率 \\
    \midrule
    良好   & xx & xx & xx\% \\
    中等   & xx & xx & xx\% \\
    恶劣   & xx & xx & xx\% \\
    \bottomrule
  \end{tabular}
\end{table}

\subsection{复现说明}

代码与数据可在 XXX 获取（详见「数据与代码可用性」声明）。


% ---------------------------------------------------------------------------
% 参考文献
% ---------------------------------------------------------------------------
\bibliographystyle{elsarticle-num-names}
\bibliography{bib/references}

\end{document}
